% =====================================================================
%  DOKUMEN TEKNIS PEATGUARD
%  Sistem Pemantauan & Peringatan Dini Kebakaran Lahan Gambut
%  Kompilasi: pdflatex dokumen-teknis-peatguard.tex  (jalankan 2 kali)
% =====================================================================
\documentclass[11pt,a4paper]{article}

% ---------- Paket dasar ----------
\usepackage[utf8]{inputenc}
\usepackage[T1]{fontenc}
\usepackage[indonesian]{babel}
\usepackage{lmodern}
\usepackage[margin=2.5cm]{geometry}
\usepackage{microtype}
\usepackage{setspace}
\setstretch{1.15}
\setlength{\parskip}{0.5em}
\setlength{\parindent}{0pt}

% ---------- Tabel & daftar ----------
\usepackage{booktabs}
\usepackage{tabularx}
\usepackage{longtable}
\usepackage{array}
\usepackage{enumitem}
\setlist{itemsep=2pt, topsep=4pt}
\newcolumntype{L}{>{\raggedright\arraybackslash}X}

% ---------- Gambar & diagram ----------
\usepackage{float}
\usepackage{caption}
\captionsetup{font=small, labelfont=bf}
\usepackage{xcolor}
\usepackage{tikz}
\usetikzlibrary{positioning, arrows.meta, shapes.multipart, shapes.geometric,
                calc, fit, backgrounds}

% ---------- Kode ----------
\usepackage{listings}
\lstset{
  basicstyle=\ttfamily\footnotesize,
  frame=single,
  rulecolor=\color{black},
  breaklines=true,
  columns=fullflexible,
  keepspaces=true,
  showstringspaces=false,
  xleftmargin=4pt, xrightmargin=4pt,
  aboveskip=8pt, belowskip=8pt,
  literate={°}{{\textdegree}}1 {·}{{$\cdot$}}1
}

% ---------- Header / footer ----------
\usepackage{fancyhdr}
\pagestyle{fancy}
\fancyhf{}
\lhead{\small Dokumen Teknis PeatGuard}
\rhead{\small Versi 1.0}
\cfoot{\small \thepage}
\renewcommand{\headrulewidth}{0.4pt}

\usepackage[hidelinks]{hyperref}

% =====================================================================
%  Gaya UML (oranye) -- tabel & gambar lain hitam-putih
% =====================================================================
\definecolor{umlfill}{RGB}{255,236,214}   % oranye sangat muda (isi)
\definecolor{umldraw}{RGB}{204,102,0}     % oranye tua (garis)

\tikzset{
  >={Stealth[length=2.4mm]},
  % kotak umum UML
  umlbox/.style={draw=umldraw, fill=umlfill, thick, rounded corners=2pt,
                 align=center, font=\small, inner sep=5pt},
  % kelas UML 3 bagian
  umlclass/.style={draw=umldraw, fill=umlfill, thick, rectangle split,
                   rectangle split parts=3, rectangle split part align={center,left,left},
                   font=\footnotesize, inner sep=4pt, text width=#1},
  umlclass/.default=5.2cm,
  % entitas ERD 2 bagian
  erd/.style={draw=umldraw, fill=umlfill, thick, rectangle split,
              rectangle split parts=2, rectangle split part align={center,left},
              font=\footnotesize, inner sep=4pt, text width=#1},
  erd/.default=5.4cm,
  % use case
  usecase/.style={draw=umldraw, fill=umlfill, thick, ellipse, align=center,
                  font=\footnotesize, minimum width=5.4cm, minimum height=0.8cm,
                  inner sep=1pt},
  % aktivitas
  action/.style={draw=umldraw, fill=umlfill, thick, rounded corners=7pt,
                 align=center, font=\small, minimum width=5.6cm,
                 minimum height=0.8cm, text width=5.4cm},
  decision/.style={draw=umldraw, fill=umlfill, thick, diamond, aspect=2.6,
                   align=center, font=\footnotesize, inner sep=1pt},
  startnode/.style={circle, fill=umldraw, minimum size=4mm, inner sep=0},
  endnode/.style={circle, draw=umldraw, thick, minimum size=5.5mm, inner sep=0,
                  path picture={\fill[umldraw] (path picture bounding box.center) circle (1.7mm);}},
  % state
  state/.style={draw=umldraw, fill=umlfill, thick, rounded corners=8pt,
                align=center, font=\small, minimum width=3.4cm, minimum height=1.2cm},
  % node deployment
  devnode/.style={draw=umldraw, fill=umlfill, thick, align=center, font=\footnotesize,
                  inner sep=5pt, text width=#1},
  devnode/.default=4.2cm,
  % panah
  flow/.style={->, thick, draw=umldraw},
  dep/.style={->, thick, dashed, draw=umldraw},
  assoc/.style={thick, draw=umldraw},
  lbl/.style={font=\scriptsize, fill=white, inner sep=1.5pt, align=center},
  % sequence diagram
  lifehead/.style={draw=umldraw, fill=umlfill, thick, font=\footnotesize,
                   align=center, minimum width=2.5cm, minimum height=0.9cm},
  msg/.style={->, thick, draw=umldraw},
  ret/.style={->, dashed, draw=umldraw},
  msglbl/.style={font=\scriptsize, above, inner sep=1.5pt, align=center},
}

% gambar aktor (stick figure)
\newcommand{\actor}[3]{% posisi, nama node, label
  \begin{scope}[shift={#1}, draw=umldraw, thick]
    \draw (0,0.55) circle (0.2);
    \draw (0,0.35) -- (0,-0.15);
    \draw (-0.35,0.2) -- (0.35,0.2);
    \draw (0,-0.15) -- (-0.28,-0.6);
    \draw (0,-0.15) -- (0.28,-0.6);
    \node[font=\footnotesize, align=center, below] at (0,-0.65) {#3};
    \coordinate (#2) at (0,0.1);
  \end{scope}
}

% kotak penjelasan
\usepackage[most]{tcolorbox}
\newtcolorbox{catatan}[1][Catatan]{
  enhanced, colback=white, colframe=black, boxrule=0.6pt, arc=0pt,
  left=6pt, right=6pt, top=4pt, bottom=4pt,
  fonttitle=\bfseries\small, coltitle=black, colbacktitle=white,
  attach title to upper, after title={\par\smallskip}, title={#1}}

\newcommand{\kode}[1]{\texttt{#1}}

% =====================================================================
\begin{document}
% =====================================================================

% ---------- Halaman judul ----------
\begin{titlepage}
  \centering
  \vspace*{2.5cm}
  {\LARGE\bfseries DOKUMEN TEKNIS SISTEM\par}
  \vspace{0.4cm}
  {\Huge\bfseries PeatGuard\par}
  \vspace{0.4cm}
  {\Large Sistem Pemantauan dan Peringatan Dini\\Kebakaran Lahan Gambut Berbasis IoT\par}
  \vspace{1.5cm}
  \rule{0.6\linewidth}{0.6pt}\par
  \vspace{0.4cm}
  {\large Bahan Pemaparan Teknis\par}
  \vspace{0.2cm}
  {\normalsize Mencakup: \textit{Backend} (\kode{be-peat}) dan \textit{Dashboard} (\kode{dashboard-peat})\par}
  \rule{0.6\linewidth}{0.6pt}\par
  \vfill
  \begin{tabular}{ll}
    Versi dokumen & 1.0 \\
    Tanggal       & 23 November 2026 \\
    Status        & Draf untuk pemaparan \\
  \end{tabular}
  \vspace*{1.5cm}
\end{titlepage}

\tableofcontents
\newpage
\listoffigures
\listoftables
\newpage

% =====================================================================
\section{Pendahuluan}
% =====================================================================

\subsection{Latar Belakang}
Kebakaran lahan gambut berbeda dengan kebakaran hutan biasa. Api di lahan gambut
sering kali \textbf{membara di bawah permukaan tanah} tanpa nyala yang terlihat,
sehingga sulit dideteksi dengan mata maupun citra satelit sampai kondisinya sudah parah.
Padahal, tanda-tandanya bisa dibaca lebih awal dari \textbf{kondisi tanah}:
tanah yang makin kering, suhu tanah yang naik, dan keasaman (pH) yang berubah.

\textbf{PeatGuard} dibuat untuk membaca tanda-tanda tersebut secara terus-menerus
(\textit{real-time}) melalui sensor di lapangan, lalu menampilkannya dalam
\textit{dashboard} web yang mudah dipahami oleh petugas pemantau.

\subsection{Tujuan Sistem}
\begin{enumerate}
  \item Memantau \textbf{kelembapan tanah}, \textbf{pH tanah}, dan \textbf{suhu tanah}
        dari beberapa titik pantau secara langsung.
  \item Menentukan \textbf{status risiko} setiap titik secara otomatis:
        \textit{Aman (normal)}, \textit{Waspada (warning)}, atau \textit{Kritis (critical)}.
  \item Memberikan \textbf{peringatan dini} ketika ada titik yang masuk kondisi berbahaya.
  \item Menyimpan \textbf{riwayat data} untuk analisis tren dan pembuatan laporan.
\end{enumerate}

\subsection{Ruang Lingkup Dokumen}
Dokumen ini menjelaskan cara kerja sistem dari sisi teknis dengan bahasa sederhana,
meliputi arsitektur, diagram UML, rancangan basis data, aturan penentuan status,
spesifikasi API, cara menjalankan, serta catatan perbaikan. Dokumen disusun
berdasarkan kode sumber pada dua repositori:

\begin{table}[H]
\centering
\caption{Repositori yang dibahas}
\label{tab:repo}
\begin{tabularx}{\linewidth}{@{}lL@{}}
\toprule
\textbf{Repositori} & \textbf{Isi} \\
\midrule
\kode{be-peat}        & \textit{Backend} API (Python, FastAPI) untuk fitur \textbf{pencarian} dan \textbf{notifikasi}. \\
\kode{dashboard-peat} & \textit{Dashboard} web (React Router), skrip \textit{bridge} MQTT ke Supabase, dan skema basis data. \\
\bottomrule
\end{tabularx}
\end{table}

\subsection{Istilah yang Perlu Diketahui}
\begin{table}[H]
\centering
\caption{Daftar istilah}
\label{tab:istilah}
\begin{tabularx}{\linewidth}{@{}lL@{}}
\toprule
\textbf{Istilah} & \textbf{Penjelasan sederhana} \\
\midrule
IoT & \textit{Internet of Things}; perangkat fisik (di sini: sensor tanah) yang bisa mengirim data lewat internet. \\
Telemetri & Data hasil pengukuran sensor yang dikirim dari jarak jauh. \\
MQTT & Protokol pengiriman pesan yang ringan, cocok untuk sensor. Bekerja seperti \textit{grup chat}: pengirim \textit{publish} ke sebuah \textbf{topik}, penerima \textit{subscribe} ke topik itu. \\
\textit{Broker} MQTT & Server ``kantor pos'' yang menerima pesan dari sensor lalu meneruskannya ke semua pihak yang berlangganan. \\
\textit{Bridge Worker} & Program kecil yang berjalan terus di latar belakang untuk memindahkan data dari MQTT ke basis data. \\
Supabase & Layanan basis data PostgreSQL di \textit{cloud} yang punya fitur \textit{Realtime} (memberi kabar ke \textit{dashboard} setiap ada data baru). \\
API & Pintu komunikasi antarprogram. \textit{Dashboard} ``bertanya'' ke API, API ``menjawab'' dalam format JSON. \\
REST & Gaya rancangan API memakai alamat (URL) dan metode HTTP seperti GET dan POST. \\
MVC & \textit{Model--View--Controller}; cara membagi kode menjadi bagian data (\textit{Model}), pintu masuk (\textit{View}), dan logika (\textit{Controller}). \\
UML & \textit{Unified Modeling Language}; standar gambar/diagram untuk menjelaskan rancangan perangkat lunak. \\
CME & \textit{Critical Moisture of Extinction}; batas kelembapan di mana bara api gambut tidak bisa lagi menjalar. \\
\bottomrule
\end{tabularx}
\end{table}

% =====================================================================
\section{Gambaran Umum Sistem}
% =====================================================================

Secara sederhana, PeatGuard bekerja dalam lima langkah:

\begin{enumerate}
  \item \textbf{Mengukur.} Sensor di tiga titik pantau mengukur kelembapan, pH, dan suhu tanah.
  \item \textbf{Mengirim.} Setiap hasil ukur dikirim ke \textit{broker} MQTT, misalnya ke topik
        \kode{sensor/soil/temp2} (suhu di titik 2).
  \item \textbf{Menilai.} \textit{Bridge worker} dan \textit{dashboard} menerima pesan itu,
        lalu menilai statusnya (Aman / Waspada / Kritis) berdasarkan ambang batas ilmiah.
  \item \textbf{Menyimpan.} Data disimpan di basis data Supabase sebagai riwayat.
        Jika statusnya kritis, dibuatkan catatan notifikasi.
  \item \textbf{Menampilkan.} \textit{Dashboard} menampilkan angka terbaru, peta, grafik tren,
        dan daftar peringatan kepada petugas.
\end{enumerate}

\begin{catatan}[Analogi singkat]
Bayangkan sensor sebagai \textbf{wartawan di lapangan}, \textit{broker} MQTT sebagai
\textbf{kantor redaksi} yang menerima dan menyebarkan berita, basis data sebagai
\textbf{arsip}, dan \textit{dashboard} sebagai \textbf{layar televisi} di ruang kendali.
Begitu ada berita baru, layar langsung ikut berubah.
\end{catatan}

\subsection{Titik Pantau}
Saat ini sistem dikonfigurasi untuk tiga titik pantau (\textit{node}):

\begin{table}[H]
\centering
\caption{Daftar titik pantau}
\label{tab:titik}
\begin{tabular}{@{}cllcc@{}}
\toprule
\textbf{Node} & \textbf{ID} & \textbf{Nama} & \textbf{Lintang} & \textbf{Bujur} \\
\midrule
1 & PG-001 & Parongpong & $-6{,}8303$ & $107{,}5867$ \\
2 & PG-002 & Baleendah  & $-6{,}8320$ & $107{,}5910$ \\
3 & PG-003 & Rancaupas  & $-6{,}8285$ & $107{,}5815$ \\
\bottomrule
\end{tabular}
\end{table}

Koordinat di atas adalah nilai awal; jika sensor mengirim \kode{latitude} dan
\kode{longitude}, posisi di peta akan diperbarui otomatis.

% =====================================================================
\section{Arsitektur Sistem}
% =====================================================================

Gambar~\ref{fig:arsitektur} menunjukkan komponen utama dan bagaimana data mengalir di antaranya.

\begin{figure}[H]
\centering
\begin{tikzpicture}[node distance=1.2cm]
  \node[umlbox, text width=3.6cm] (sensor) at (0,0)
     {\textbf{Perangkat Sensor}\\ 3 titik pantau\\ \scriptsize kelembapan, pH, suhu};
  \node[umlbox, text width=3.6cm] (broker) at (5.2,0)
     {\textbf{Broker MQTT}\\ \scriptsize via WebSocket (wss)\\ \scriptsize topik \kode{sensor/soil/\#}};
  \node[umlbox, text width=3.6cm] (bridge) at (10.4,0)
     {\textbf{Bridge Worker}\\ \scriptsize Node.js\\ \scriptsize \kode{mqtt\_supabase\_bridge.js}};
  \node[umlbox, text width=3.6cm] (db) at (10.4,-3.6)
     {\textbf{Supabase}\\ \scriptsize PostgreSQL + Realtime\\ \scriptsize 3 tabel};
  \node[umlbox, text width=3.6cm] (dash) at (5.2,-3.6)
     {\textbf{Dashboard Web}\\ \scriptsize React Router + Leaflet\\ \scriptsize 6 halaman};
  \node[umlbox, text width=3.6cm] (api) at (0,-3.6)
     {\textbf{Backend API}\\ \scriptsize Python FastAPI (MVC)\\ \scriptsize pencarian \& notifikasi};
  \actor{(5.2,-6.6)}{user}{Petugas / Pengguna}

  \draw[flow] (sensor) -- node[lbl]{publish} (broker);
  \draw[flow] (broker) -- node[lbl]{subscribe} (bridge);
  \draw[flow] (bridge) -- node[lbl, right=2pt, fill=none]{simpan log,\\perbarui status,\\notifikasi} (db);
  \draw[flow] (db) -- node[lbl]{data riwayat\\+ realtime} (dash);
  \draw[flow] (broker) -- node[lbl]{subscribe\\langsung} (dash);
  \draw[flow, <->] (dash) -- node[lbl]{HTTP\\REST/JSON} (api);
  \draw[flow, <->] ($(user)+(0,0.55)$) -- node[lbl, right=2pt, fill=none]{browser} (dash);
\end{tikzpicture}
\caption{Arsitektur umum PeatGuard}
\label{fig:arsitektur}
\end{figure}

\subsection{Penjelasan Komponen}
\begin{table}[H]
\centering
\caption{Peran tiap komponen}
\label{tab:komponen}
\begin{tabularx}{\linewidth}{@{}lL@{}}
\toprule
\textbf{Komponen} & \textbf{Peran} \\
\midrule
Perangkat Sensor & Mengukur kondisi tanah dan mengirim angka ke \textit{broker} MQTT. Satu titik pantau mengirim 3 topik (kelembapan, pH, suhu), jadi total ada 9 topik. \\
Broker MQTT & Menerima pesan dari sensor dan meneruskannya ke \textit{bridge worker} dan \textit{dashboard}. Diakses lewat WebSocket aman (\kode{wss://}). \\
Bridge Worker & Berjalan 24 jam. Menyimpan setiap pesan ke Supabase, menghitung status risiko, dan membuat notifikasi saat kondisi kritis. Ini membuat data tetap tersimpan walaupun tidak ada yang membuka \textit{dashboard}. \\
Supabase & Basis data utama: menyimpan data titik sensor, riwayat pengukuran, dan notifikasi. Fitur \textit{Realtime} memberi tahu \textit{dashboard} setiap ada baris baru. \\
Dashboard Web & Antarmuka bagi petugas: ringkasan, peta, kondisi lahan, peringatan, riwayat \& grafik, serta laporan CSV. \\
Backend API & Melayani fitur \textbf{pencarian} (sensor dan halaman) dan \textbf{notifikasi} (daftar, tandai dibaca). Dirancang dengan pola MVC. \\
\bottomrule
\end{tabularx}
\end{table}

\subsection{Teknologi yang Digunakan}
\begin{table}[H]
\centering
\caption{Daftar teknologi}
\label{tab:teknologi}
\begin{tabularx}{\linewidth}{@{}llL@{}}
\toprule
\textbf{Bagian} & \textbf{Teknologi} & \textbf{Kegunaan} \\
\midrule
Backend      & Python, FastAPI $\geq$ 0.100      & Kerangka kerja (\textit{framework}) API. \\
             & Uvicorn $\geq$ 0.20               & Server untuk menjalankan FastAPI (port 8000). \\
             & Pydantic $\geq$ 2.0               & Validasi \& bentuk data (Model). \\
\midrule
Frontend     & React 19, React Router 8          & Membangun halaman \textit{dashboard} (dengan \textit{server-side rendering}). \\
             & Tailwind CSS                      & Tampilan/gaya antarmuka. \\
             & Leaflet                           & Peta interaktif sebaran sensor. \\
             & mqtt.js                           & Klien MQTT di \textit{browser} dan di \textit{bridge}. \\
             & supabase-js                       & Klien basis data \& \textit{realtime}. \\
             & Vite, TypeScript                  & Alat \textit{build} dan pemeriksaan tipe. \\
\midrule
Data         & Supabase (PostgreSQL)             & Penyimpanan data, keamanan baris (RLS), \textit{realtime}. \\
\midrule
Infrastruktur& Docker (Node 24 Alpine)           & Mengemas \textit{dashboard} agar mudah dipasang (port 3000). \\
\bottomrule
\end{tabularx}
\end{table}

% =====================================================================
\section{Pemodelan Sistem dengan UML}
% =====================================================================

Bagian ini menggambarkan sistem dari beberapa sudut pandang memakai diagram UML.
Setiap diagram dilengkapi penjelasan singkat cara membacanya.

% ---------------------------------------------------------------------
\subsection{Use Case Diagram}
% ---------------------------------------------------------------------
\textit{Use case diagram} menjawab pertanyaan: \textbf{siapa melakukan apa} terhadap sistem.
Ada tiga aktor: \textbf{Petugas} (manusia), \textbf{Perangkat Sensor}, dan
\textbf{Bridge Worker} (program otomatis).

\begin{figure}[H]
\centering
\begin{tikzpicture}
  % batas sistem
  \draw[umldraw, thick] (2.0,0.9) rectangle (10.6,-10.4);
  \node[font=\small\bfseries, anchor=north] at (6.3,0.85) {Sistem PeatGuard};

  \node[usecase] (uc1)  at (6.3,0.0)   {Melihat ringkasan dashboard};
  \node[usecase] (uc2)  at (6.3,-0.95) {Melihat peta sebaran sensor};
  \node[usecase] (uc3)  at (6.3,-1.9)  {Melihat kondisi fisik \& kimia lahan};
  \node[usecase] (uc4)  at (6.3,-2.85) {Melihat peringatan dini};
  \node[usecase] (uc5)  at (6.3,-3.8)  {Melihat riwayat \& grafik tren};
  \node[usecase] (uc6)  at (6.3,-4.75) {Mengunduh laporan (CSV)};
  \node[usecase] (uc7)  at (6.3,-5.7)  {Mencari sensor / halaman};
  \node[usecase] (uc8)  at (6.3,-6.65) {Membaca \& menandai notifikasi};
  \node[usecase] (uc9)  at (6.3,-7.6)  {Mengirim data telemetri};
  \node[usecase] (uc10) at (6.3,-8.55) {Menyimpan data ke basis data};
  \node[usecase] (uc11) at (6.3,-9.7)  {Menentukan status risiko};

  \actor{(0,-3.4)}{petugas}{Petugas}
  \actor{(12.6,-6.9)}{sensor}{Perangkat\\Sensor}
  \actor{(12.6,-9.9)}{bridge}{Bridge\\Worker}

  \foreach \u in {uc1,uc2,uc3,uc4,uc5,uc6,uc7,uc8}
     \draw[assoc] (petugas) -- (\u.west);
  \draw[assoc] (sensor) -- (uc9.east);
  \draw[assoc] (bridge) -- (uc10.east);
  \draw[dep] (uc10) -- node[lbl, right=1pt, fill=none]{\guillemotleft include\guillemotright} (uc11);
\end{tikzpicture}
\caption{Use case diagram PeatGuard}
\label{fig:usecase}
\end{figure}

\begin{table}[H]
\centering
\caption{Penjelasan use case}
\label{tab:usecase}
\begin{tabularx}{\linewidth}{@{}lL@{}}
\toprule
\textbf{Use case} & \textbf{Penjelasan} \\
\midrule
Melihat ringkasan dashboard & Melihat kartu angka terbaru tiap titik, status koneksi MQTT \& Supabase, dan 9 topik telemetri. \\
Melihat peta sebaran sensor & Melihat posisi titik pantau di peta Leaflet beserta warna statusnya. \\
Melihat kondisi lahan & Melihat rata-rata kelembapan, pH, suhu, dan tabel rincian per titik. \\
Melihat peringatan dini & Melihat titik yang berstatus Waspada/Kritis beserta alasannya. \\
Melihat riwayat \& tren & Melihat grafik naik-turun nilai sensor dan tabel log (bisa dicari \& diurutkan). \\
Mengunduh laporan & Mengekspor riwayat data menjadi berkas CSV. \\
Mencari sensor / halaman & Mengetik kata kunci (mis. ``waspada'', ``PG-002'', ``peta'') pada kolom pencarian. \\
Membaca notifikasi & Membuka daftar notifikasi, menandai satu atau semua sebagai sudah dibaca. \\
Mengirim data telemetri & Sensor mengirim angka hasil ukur ke \textit{broker} MQTT. \\
Menyimpan data & \textit{Bridge} menyimpan data ke Supabase dan \textbf{selalu} menyertakan penentuan status risiko (\guillemotleft include\guillemotright). \\
\bottomrule
\end{tabularx}
\end{table}

% ---------------------------------------------------------------------
\subsection{Activity Diagram: Alur Data Sensor}
% ---------------------------------------------------------------------
\textit{Activity diagram} menggambarkan \textbf{urutan langkah} dari satu proses,
termasuk percabangan keputusan. Gambar~\ref{fig:activity} memperlihatkan perjalanan
satu pesan sensor sampai tampil di \textit{dashboard}.

\begin{figure}[H]
\centering
\begin{tikzpicture}[node distance=0.55cm]
  \node[startnode] (s) {};
  \node[action, below=of s] (a1) {Sensor mengukur kelembapan, pH, atau suhu};
  \node[action, below=of a1] (a2) {Sensor mengirim angka ke topik MQTT\\ \scriptsize contoh: \kode{sensor/soil/moisture1}};
  \node[action, below=of a2] (a3) {Bridge Worker menerima pesan};
  \node[decision, below=of a3] (d1) {Angka \& nomor\\titik valid?};
  \node[action, below=0.7cm of d1] (a4) {Perbarui nilai terakhir titik tersebut\\(cache)};
  \node[action, below=of a4] (a5) {Hitung status risiko \& alasannya};
  \node[action, below=of a5] (a6) {Simpan ke \kode{sensor\_logs}\\ dan perbarui tabel \kode{sensors}};
  \node[decision, below=of a6] (d2) {Status\\kritis?};
  \node[action, right=1.2cm of d2, minimum width=3.6cm, text width=3.4cm] (a7) {Buat catatan di \kode{notifications}};
  \coordinate[below=0.6cm of d2] (m);
  \node[action, below=0.5cm of m] (a8) {Dashboard menerima data baru (realtime) dan memperbarui tampilan};
  \node[endnode, below=of a8] (e) {};

  \node[action, left=1.2cm of d1, minimum width=2.6cm, text width=2.4cm] (x) {Abaikan pesan};
  \node[endnode, below=0.5cm of x] (ex) {};

  \draw[flow] (s) -- (a1);
  \draw[flow] (a1) -- (a2);
  \draw[flow] (a2) -- (a3);
  \draw[flow] (a3) -- (d1);
  \draw[flow] (d1) -- node[lbl, right]{ya} (a4);
  \draw[flow] (d1) -- node[lbl, above]{tidak} (x);
  \draw[flow] (x) -- (ex);
  \draw[flow] (a4) -- (a5);
  \draw[flow] (a5) -- (a6);
  \draw[flow] (a6) -- (d2);
  \draw[flow] (d2) -- node[lbl, above]{ya} (a7);
  \draw[thick, draw=umldraw] (d2) -- node[lbl, right]{tidak} (m);
  \draw[thick, draw=umldraw] (a7) |- (m);
  \draw[flow] (m) -- (a8);
  \draw[flow] (a8) -- (e);
\end{tikzpicture}
\caption{Activity diagram alur data dari sensor ke dashboard}
\label{fig:activity}
\end{figure}

% ---------------------------------------------------------------------
\subsection{Sequence Diagram: Pengiriman Telemetri}
% ---------------------------------------------------------------------
\textit{Sequence diagram} menunjukkan \textbf{siapa berbicara kepada siapa, dan dalam urutan apa}.
Waktu berjalan dari atas ke bawah. Panah penuh adalah permintaan/pesan, panah putus-putus adalah balasan
atau pemberitahuan.

\begin{figure}[H]
\centering
\begin{tikzpicture}
  \foreach \x/\n/\t in {0/S/Sensor, 3.1/B/Broker\\MQTT, 6.2/W/Bridge\\Worker, 9.3/D/Supabase, 12.4/F/Dashboard}{
    \node[lifehead] (\n) at (\x,0) {\t};
    \draw[dashed, draw=umldraw!70] (\n.south) -- (\x,-12.2);
  }
  % aktivasi
  \fill[umlfill, draw=umldraw] (6.1,-1.3) rectangle (6.3,-8.6);
  \fill[umlfill, draw=umldraw] (12.3,-2.1) rectangle (12.5,-4.4);

  \draw[msg] (0,-1.0) -- node[msglbl]{1: publish(\kode{temp2}, 33.7)} (3.1,-1.0);
  \draw[msg] (3.1,-1.5) -- node[msglbl]{2: teruskan pesan} (6.1,-1.5);
  \draw[msg] (3.1,-2.3) -- node[msglbl, pos=0.78]{2a: teruskan pesan} (12.3,-2.3);
  % self call dashboard (label di kiri)
  \draw[msg] (12.3,-2.9) -- ++(-0.6,0) -- ++(0,-0.4) -- ++(0.6,0);
  \node[font=\scriptsize, anchor=east, align=right] at (11.6,-3.1) {2b: hitung status, perbarui kartu};
  \draw[msg] (12.3,-4.1) -- node[msglbl]{2c: simpan log (dari browser)} (9.3,-4.1);

  % self call bridge
  \draw[msg] (6.3,-4.9) -- ++(0.6,0) -- ++(0,-0.4) -- ++(-0.6,0);
  \node[font=\scriptsize, anchor=west] at (7.0,-5.1) {3: evaluateStatus()};
  \draw[msg] (6.3,-6.0) -- node[msglbl]{4: insert \kode{sensor\_logs}} (9.3,-6.0);
  \draw[msg] (6.3,-6.8) -- node[msglbl]{5: update \kode{sensors}} (9.3,-6.8);

  % frame opt
  \draw[umldraw, thick] (5.5,-7.2) rectangle (10.2,-8.9);
  \node[draw=umldraw, fill=white, font=\scriptsize\bfseries, anchor=north west, inner sep=2pt] at (5.5,-7.2) {opt};
  \node[font=\scriptsize, anchor=north west] at (6.4,-7.25) {[status = kritis]};
  \draw[msg] (6.3,-8.3) -- node[msglbl]{6: insert \kode{notifications}} (9.3,-8.3);

  \fill[umlfill, draw=umldraw] (12.3,-9.5) rectangle (12.5,-11.2);
  \draw[ret] (9.3,-9.7) -- node[msglbl]{7: event INSERT} (12.3,-9.7);
  \draw[msg] (12.3,-10.3) -- ++(-0.6,0) -- ++(0,-0.4) -- ++(0.6,0);
  \node[font=\scriptsize, anchor=east, align=right] at (11.6,-10.5) {8: perbarui tabel \& grafik riwayat};
\end{tikzpicture}
\caption{Sequence diagram pengiriman dan penyimpanan data telemetri}
\label{fig:seq-telemetri}
\end{figure}

\textbf{Cara membaca:} sensor mengirim satu angka (langkah 1). \textit{Broker} meneruskannya ke
\textit{bridge} (2) \textbf{dan} langsung ke \textit{dashboard} yang sedang dibuka (2a), sehingga angka di layar
berubah seketika (2b). \textit{Bridge} menghitung status (3), menyimpan ke basis data (4--5), dan membuat
notifikasi jika kritis (6). Supabase kemudian memberi kabar ke \textit{dashboard} bahwa ada baris riwayat baru (7--8).

\begin{catatan}[Perhatian]
Pada langkah 2c, \textit{dashboard} yang terbuka di \textit{browser} juga ikut menyimpan log ke Supabase.
Akibatnya, satu pesan sensor bisa tercatat \textbf{lebih dari satu kali} (sekali oleh \textit{bridge},
dan sekali oleh setiap \textit{browser} yang terbuka). Lihat Bagian~\ref{sec:temuan}.
\end{catatan}

% ---------------------------------------------------------------------
\subsection{Sequence Diagram: Pencarian dan Notifikasi}
% ---------------------------------------------------------------------
Fitur pencarian dan notifikasi di \textit{header dashboard} dilayani oleh \textit{backend} FastAPI.

\begin{figure}[H]
\centering
\begin{tikzpicture}
  \actor{(0,0.1)}{p}{}
  \node[font=\footnotesize] at (0,-0.75) {Petugas};
  \foreach \x/\n/\t in {3.4/F/Dashboard\\(Header), 6.8/V/View\\(Router), 10.2/C/Controller, 13.2/M/Model}{
    \node[lifehead, minimum width=2.3cm] (\n) at (\x,0) {\t};
  }
  \foreach \x in {0,3.4,6.8,10.2,13.2}
    \draw[dashed, draw=umldraw!70] (\x,-1.0) -- (\x,-11.4);

  \node[font=\scriptsize\bfseries, anchor=west] at (-0.4,-1.3) {A. Pencarian};
  \draw[msg] (0,-1.8) -- node[msglbl]{1: ketik ``waspada''} (3.4,-1.8);
  \draw[msg] (3.4,-2.5) -- node[msglbl]{2: GET \kode{/api/search?q=waspada}} (6.8,-2.5);
  \draw[msg] (6.8,-3.1) -- node[msglbl]{3: search(q)} (10.2,-3.1);
  \draw[msg] (10.2,-3.7) -- node[msglbl]{4: cocokkan sensor\\\& halaman} (13.2,-3.7);
  \draw[ret] (10.2,-4.3) -- node[msglbl]{5: SearchResponse} (6.8,-4.3);
  \draw[ret] (6.8,-4.9) -- node[msglbl]{6: JSON hasil} (3.4,-4.9);
  \draw[ret] (3.4,-5.5) -- node[msglbl]{7: tampilkan daftar hasil} (0,-5.5);

  \node[font=\scriptsize\bfseries, anchor=west] at (-0.4,-6.3) {B. Tandai notifikasi dibaca};
  \draw[msg] (0,-6.8) -- node[msglbl]{8: klik notifikasi} (3.4,-6.8);
  \draw[msg] (3.4,-7.5) -- node[msglbl]{9: POST \kode{/api/notifications/\{id\}/read}} (6.8,-7.5);
  \draw[msg] (6.8,-8.1) -- node[msglbl]{10: mark\_as\_read(id)} (10.2,-8.1);
  \draw[msg] (10.2,-8.7) -- node[msglbl]{11: is\_read = true} (13.2,-8.7);
  \draw[ret] (10.2,-9.3) -- node[msglbl]{12: MarkReadResponse} (6.8,-9.3);
  \draw[ret] (6.8,-9.9) -- node[msglbl]{13: JSON} (3.4,-9.9);
  \draw[ret] (3.4,-10.6) -- node[msglbl]{14: buka halaman terkait} (0,-10.6);
\end{tikzpicture}
\caption{Sequence diagram fitur pencarian dan notifikasi}
\label{fig:seq-api}
\end{figure}

\begin{catatan}[Mekanisme cadangan (\textit{fallback})]
Jika \textit{backend} tidak menjawab dalam 2 detik, \textit{dashboard} tetap berfungsi:
pencarian dilakukan langsung di \textit{browser} dan notifikasi memakai data contoh bawaan.
Daftar notifikasi juga diperbarui otomatis setiap 10 detik (\textit{polling}).
\end{catatan}

% ---------------------------------------------------------------------
\subsection{Class Diagram: Backend (Pola MVC)}
% ---------------------------------------------------------------------
\textit{Backend} memakai pola MVC. Pembagiannya:
\begin{itemize}
  \item \textbf{View} (\kode{app/views}) -- pintu masuk HTTP; hanya menerima permintaan dan meneruskannya.
  \item \textbf{Controller} (\kode{app/controllers}) -- tempat logika: mencari, menghitung, mengubah data.
  \item \textbf{Model} (\kode{app/models}) -- bentuk data yang dipakai dan dikembalikan (Pydantic).
\end{itemize}

\begin{figure}[H]
\centering
\begin{tikzpicture}
  \node[umlclass=5.6cm, anchor=north] (sv) at (0,0)
    {\guillemotleft router\guillemotright\ \textbf{search\_view}
     \nodepart{second} prefix = "/api/search"
     \nodepart{third} + search\_endpoint(q): SearchResponse};
  \node[umlclass=5.8cm, anchor=north] (sc) at (0,-2.5)
    {\textbf{SearchController}
     \nodepart{second} + SENSORS\_DATA: List[SensorModel]\\ + PAGES\_DATA: List[dict]
     \nodepart{third} + search(query: str): SearchResponse};
  \node[umlclass=4.4cm, anchor=north] (sm) at (-4.0,-5.6)
    {\textbf{SensorModel}
     \nodepart{second} + id: str\\ + name: str\\ + region: str\\ + lat, lng: float\\ + status: str\\ + moisture: float\\ + ph: float\\ + temperature: float\\ + water\_level: float?\\ + last\_reading: str
     \nodepart{third} \textit{(tidak ada method)}};
  \node[umlclass=4.6cm, anchor=north] (sr) at (4.0,-5.6)
    {\textbf{SearchResponse}
     \nodepart{second} + query: str\\ + total\_results: int\\ + results: List[SearchResultItem]
     \nodepart{third} \textit{(tidak ada method)}};
  \node[umlclass=4.6cm, anchor=north] (si) at (4.0,-9.4)
    {\textbf{SearchResultItem}
     \nodepart{second} + id: str\\ + title: str\\ + category: str\\ + subtitle: str\\ + status: str?\\ + url: str
     \nodepart{third} \textit{(tidak ada method)}};

  \draw[dep] (sv) -- node[lbl, right=2pt, fill=none]{memanggil} (sc);
  \draw[assoc, {Diamond[open, length=3mm]}-] ($(sc.south)+(-2.0,0)$) -- ++(0,-0.7) -| node[lbl, pos=0.8, right=2pt]{1..*} (sm.north);
  \draw[dep] ($(sc.south)+(2.0,0)$) -- ++(0,-0.7) -| node[lbl, pos=0.8, left=2pt]{membuat} (sr.north);
  \draw[assoc, {Diamond[length=3mm]}-] (sr.south) -- node[lbl, right=2pt, fill=none]{0..*} (si.north);
\end{tikzpicture}
\caption{Class diagram fitur pencarian (backend)}
\label{fig:class-search}
\end{figure}

\begin{figure}[H]
\centering
\begin{tikzpicture}
  \node[umlclass=7.4cm, anchor=north] (nv) at (-2.2,0)
    {\guillemotleft router\guillemotright\ \textbf{notification\_view}
     \nodepart{second} prefix = "/api/notifications"
     \nodepart{third} + get\_notifications(): NotificationListResponse\\
                      + mark\_as\_read(notif\_id): MarkReadResponse\\
                      + mark\_all\_as\_read(): MarkReadResponse\\
                      + create\_notification(req): NotificationModel};
  \node[umlclass=4.0cm, anchor=north] (cr) at (5.0,0)
    {\textbf{CreateNotification-\\Request}
     \nodepart{second} + title: str\\ + message: str\\ + type: str = "info"\\ + sensor\_id: str?\\ + url: str = "/"
     \nodepart{third} \textit{(tidak ada method)}};
  \node[umlclass=7.4cm, anchor=north] (nc) at (-2.2,-3.9)
    {\textbf{NotificationController}
     \nodepart{second} - \_notifications: List[NotificationModel] \textit{(di memori)}
     \nodepart{third} + get\_all(): NotificationListResponse\\
                      + mark\_as\_read(notif\_id: str): MarkReadResponse\\
                      + mark\_all\_as\_read(): MarkReadResponse\\
                      + add\_notification(...): NotificationModel};
  \node[umlclass=3.9cm, anchor=north] (nm) at (-4.6,-8.2)
    {\textbf{NotificationModel}
     \nodepart{second} + id: str\\ + title: str\\ + message: str\\ + type: str\\ + timestamp: str\\ + created\_at: datetime\\ + is\_read: bool\\ + sensor\_id: str?\\ + url: str
     \nodepart{third} \textit{(tidak ada method)}};
  \node[umlclass=4.2cm, anchor=north] (nl) at (0.6,-8.2)
    {\textbf{NotificationListResponse}
     \nodepart{second} + unread\_count: int\\ + total\_count: int\\ + notifications: List[...]
     \nodepart{third} \textit{(tidak ada method)}};
  \node[umlclass=3.6cm, anchor=north] (mr) at (5.2,-8.2)
    {\textbf{MarkReadResponse}
     \nodepart{second} + success: bool\\ + message: str\\ + unread\_count: int
     \nodepart{third} \textit{(tidak ada method)}};

  \draw[dep] (nv) -- node[lbl, right=2pt, fill=none]{memanggil} (nc);
  \draw[dep] (nv.east) -- node[lbl, above=2pt, fill=none]{memakai} (cr.west |- nv.east);
  \draw[assoc, {Diamond[open, length=3mm]}-] (nc.south -| nm.north) -- node[lbl, right=2pt, fill=none]{0..*} (nm.north);
  \draw[dep] (nc.south -| nl.north) -- node[lbl, right=2pt, fill=none]{membuat} (nl.north);
  \draw[dep] (nc.east |- 0,-6.0) -| node[lbl, pos=0.3, above=2pt, fill=none]{membuat} (mr.north);
  \draw[assoc, {Diamond[length=3mm]}-] (nl.west) -- node[lbl, below=1pt, fill=none]{0..*} (nm.east |- nl.west);
\end{tikzpicture}
\caption{Class diagram fitur notifikasi (backend)}
\label{fig:class-notif}
\end{figure}

\textbf{Arti simbol:} panah putus-putus berarti ``memanggil/membuat'';
belah ketupat kosong berarti ``memiliki kumpulan'' (agregasi);
belah ketupat penuh berarti ``terdiri dari'' (komposisi);
angka seperti \kode{0..*} berarti ``nol atau lebih''.

% ---------------------------------------------------------------------
\subsection{Component Diagram: Dashboard (Frontend)}
% ---------------------------------------------------------------------
Gambar~\ref{fig:komponen} menunjukkan susunan bagian-bagian kode \textit{dashboard}.
Kunci utamanya adalah \textbf{SensorContext}: satu ``pusat data'' yang dipakai bersama oleh semua halaman,
sehingga koneksi MQTT dan Supabase cukup dibuat sekali.

\begin{figure}[H]
\centering
\begin{tikzpicture}
  \node[umlbox, text width=5cm] (root) at (0,0) {\guillemotleft component\guillemotright\\\textbf{root.jsx} (App)};
  \node[umlbox, text width=5cm] (ctx) at (0,-1.8) {\guillemotleft component\guillemotright\\\textbf{SensorContext}\\\scriptsize (SensorProvider, useSensorData)};

  \node[umlbox, text width=3.0cm] (lib) at (-4.4,-4.2) {\guillemotleft module\guillemotright\\\textbf{lib/supabase.js}\\\scriptsize ambil \& simpan data};
  \node[umlbox, text width=3.0cm] (data) at (0,-4.2) {\guillemotleft module\guillemotright\\\textbf{data/sensors.js}\\\scriptsize data awal \& ambang batas};
  \node[umlbox, text width=3.0cm] (mq) at (4.4,-4.2) {\guillemotleft library\guillemotright\\\textbf{mqtt.js}\\\scriptsize klien WebSocket};

  \node[umlbox, text width=3.0cm] (sb) at (-4.4,-6.3) {\guillemotleft external\guillemotright\\\textbf{Supabase}};
  \node[umlbox, text width=3.0cm] (br) at (4.4,-6.3) {\guillemotleft external\guillemotright\\\textbf{Broker MQTT}};

  \node[umlbox, text width=11.0cm, align=left, font=\footnotesize] (pages) at (0,-8.6)
    {\textbf{Halaman (routes)}\\[2pt]
     \kode{home.jsx}\ $\cdot$\ \kode{riwayat.jsx}\ $\cdot$\ \kode{peta-sensor.jsx}\ $\cdot$\
     \kode{kondisi-lahan.jsx}\ $\cdot$\ \kode{peringatan.jsx}\ $\cdot$\ \kode{laporan.jsx}};
  \node[umlbox, text width=11.0cm, align=left, font=\footnotesize] (comp) at (0,-10.7)
    {\textbf{Komponen UI (components)}\\[2pt]
     \kode{DashboardFrame}\ $\cdot$\ \kode{Sidebar}\ $\cdot$\ \kode{MainMenu}\ $\cdot$\
     \kode{SensorMap}\ $\cdot$\ \kode{DataTable}\ $\cdot$\ \kode{ThresholdReferenceCard}\ $\cdot$\ \kode{Footer}};
  \node[umlbox, text width=3.4cm] (api) at (0,-12.8) {\guillemotleft external\guillemotright\\\textbf{Backend FastAPI}};

  \draw[dep] (root) -- node[lbl, right=2pt, fill=none]{membungkus} (ctx);
  \draw[dep] (ctx) -- (lib);
  \draw[dep] (ctx) -- (data);
  \draw[dep] (ctx) -- (mq);
  \draw[dep] (lib) -- (sb);
  \draw[dep] (mq) -- (br);
  \draw[dep] (pages.west) -- ++(-0.6,0) |- node[lbl, pos=0.25, rotate=90]{useSensorData()} (ctx.west);
  \draw[dep] (pages) -- node[lbl, right=2pt, fill=none]{memakai} (comp);
  \draw[dep] (comp) -- node[lbl, right=2pt, fill=none]{\kode{DashboardFrame}: pencarian \& notifikasi} (api);
\end{tikzpicture}
\caption{Component diagram dashboard}
\label{fig:komponen}
\end{figure}

% ---------------------------------------------------------------------
\subsection{State Diagram: Status Titik Pantau}
% ---------------------------------------------------------------------
Setiap titik pantau selalu berada di salah satu dari tiga keadaan. Keadaan berubah otomatis
setiap kali ada nilai sensor baru. Singkatan: \textbf{K} = kelembapan (\%), \textbf{S} = suhu (\textdegree C).

\begin{figure}[H]
\centering
\begin{tikzpicture}
  \node[startnode] (st) at (-6.6,0) {};
  \node[state] (n) at (-3.6,0) {\textbf{Aman}\\\scriptsize (normal)};
  \node[state] (w) at (3.6,0) {\textbf{Waspada}\\\scriptsize (warning)};
  \node[state] (c) at (0,-4.4) {\textbf{Kritis}\\\scriptsize (critical)};

  \draw[flow] (st) -- (n);
  \draw[flow] ([yshift=4pt]n.east) -- node[lbl, above=2pt, fill=none]{K $<$ 60 atau S $\geq$ 32\\atau pH di luar 3,2--5,5} ([yshift=4pt]w.west);
  \draw[flow] ([yshift=-4pt]w.west) -- node[lbl, below=2pt, fill=none]{semua kembali normal} ([yshift=-4pt]n.east);
  \draw[flow] ([xshift=-8pt]n.south) -- node[lbl, left=3pt, fill=none]{K $<$ 40 atau\\S $\geq$ 35} ([xshift=-8pt]c.north west);
  \draw[flow] ([xshift=8pt]c.north west) -- node[lbl, right=3pt, fill=none, pos=0.4]{semua\\normal} ([xshift=8pt]n.south);
  \draw[flow] ([xshift=8pt]w.south) -- node[lbl, right=3pt, fill=none]{K $<$ 40 atau\\S $\geq$ 35} ([xshift=8pt]c.north east);
  \draw[flow] ([xshift=-8pt]c.north east) -- node[lbl, left=3pt, fill=none, pos=0.4]{K $\geq$ 40\\dan S $<$ 35} ([xshift=-8pt]w.south);
\end{tikzpicture}
\caption{State diagram status risiko titik pantau}
\label{fig:state}
\end{figure}

% ---------------------------------------------------------------------
\subsection{Deployment Diagram}
% ---------------------------------------------------------------------
\textit{Deployment diagram} menunjukkan \textbf{di mana} setiap program dijalankan.

\begin{figure}[H]
\centering
\begin{tikzpicture}
  \node[devnode=3.4cm] (dev) at (-5.0,0)
    {\guillemotleft device\guillemotright\\\textbf{Perangkat Sensor IoT}\\[2pt]\scriptsize 3 node lapangan};
  \node[devnode=3.4cm] (brk) at (0,0)
    {\guillemotleft server\guillemotright\\\textbf{Server Broker MQTT}\\[2pt]\scriptsize WebSocket TLS (wss)};
  \node[devnode=3.4cm] (cloud) at (5.0,0)
    {\guillemotleft cloud\guillemotright\\\textbf{Supabase}\\[2pt]\scriptsize PostgreSQL + Realtime};

  \node[devnode=4.8cm] (app) at (-3.0,-3.6)
    {\guillemotleft server\guillemotright\\\textbf{Server Aplikasi}\\[2pt]
     \scriptsize Docker \kode{node:24-alpine}: \kode{react-router-serve} (port 3000)\\
     \scriptsize Proses Node.js: \textit{bridge worker} (\kode{npm run bridge})};
  \node[devnode=3.8cm] (be) at (3.6,-3.6)
    {\guillemotleft server\guillemotright\\\textbf{Server Backend}\\[2pt]
     \scriptsize Python + Uvicorn, FastAPI (port 8000)\\
     \scriptsize Notifikasi di memori};
  \node[devnode=3.4cm] (br) at (0,-7.2)
    {\guillemotleft device\guillemotright\\\textbf{Browser Pengguna}\\[2pt]\scriptsize PC / laptop / ponsel};

  \draw[assoc] (dev) -- node[lbl]{MQTT} (brk);
  \draw[assoc] (brk.south) -- node[lbl]{wss} (app.north east);
  \draw[assoc] (cloud.south west) -- node[lbl, pos=0.4]{HTTPS} (app.north east);
  \draw[assoc] (br) -- node[lbl]{HTTP :3000} (app);
  \draw[assoc] (br) -- node[lbl]{HTTP :8000} (be);
  \draw[assoc] (br.west) -- (-7.4,-7.2) |- node[lbl, pos=0.25, rotate=90]{wss (langsung ke broker)} (0,1.4) -- (brk.north);
  \draw[assoc] (br.east) -- (7.4,-7.2) |- node[lbl, pos=0.25, rotate=90]{HTTPS + Realtime} (cloud.east);
\end{tikzpicture}
\caption{Deployment diagram (usulan penempatan)}
\label{fig:deploy}
\end{figure}

Garis di sisi kiri dan kanan gambar menandakan bahwa \textit{browser} pengguna juga terhubung
\textbf{langsung} ke \textit{broker} MQTT (wss) dan ke Supabase (HTTPS + \textit{Realtime}).

% =====================================================================
\section{Rancangan Basis Data}
% =====================================================================

Basis data berada di Supabase (PostgreSQL) dengan tiga tabel. Hubungan antartabel ditunjukkan pada
\textit{Entity Relationship Diagram} (ERD) berikut.

\begin{figure}[H]
\centering
\begin{tikzpicture}
  \node[erd=5.0cm, anchor=north] (s) at (0,0)
    {\textbf{sensors}
     \nodepart{second}
     \textbf{PK} id : TEXT\\
     name : TEXT\\
     region : TEXT\\
     latitude : DOUBLE\\
     longitude : DOUBLE\\
     last\_moisture : DOUBLE\\
     last\_ph : DOUBLE\\
     last\_temperature : DOUBLE\\
     status : TEXT\\
     updated\_at : TIMESTAMPTZ};
  \node[erd=5.0cm, anchor=north] (l) at (8.2,0)
    {\textbf{sensor\_logs}
     \nodepart{second}
     \textbf{PK} id : BIGINT (otomatis)\\
     created\_at : TIMESTAMPTZ\\
     \textbf{FK} sensor\_id : TEXT\\
     sensor\_name : TEXT\\
     moisture : DOUBLE\\
     ph : DOUBLE\\
     temperature : DOUBLE\\
     latitude : DOUBLE\\
     longitude : DOUBLE\\
     status : TEXT\\
     reason : TEXT};
  \node[erd=5.0cm, anchor=north] (n) at (0,-6.6)
    {\textbf{notifications}
     \nodepart{second}
     \textbf{PK} id : UUID\\
     created\_at : TIMESTAMPTZ\\
     \textbf{FK} sensor\_id : TEXT\\
     title : TEXT\\
     message : TEXT\\
     severity : TEXT\\
     is\_read : BOOLEAN};

  \draw[assoc] (s.east |- 0,-1.6) -- node[lbl, pos=0.1, above]{1} node[lbl, pos=0.9, above]{N}
               node[lbl, below=2pt, fill=none]{memiliki} (l.west |- 0,-1.6);
  \draw[assoc] (s.south) -- node[lbl, pos=0.15, right]{1} node[lbl, pos=0.85, right]{N}
               node[lbl, left=3pt, fill=none]{memiliki} (n.north);
\end{tikzpicture}
\caption{Entity Relationship Diagram (ERD)}
\label{fig:erd}
\end{figure}

\textbf{Penjelasan singkat:}
\begin{itemize}
  \item \kode{sensors} menyimpan \textbf{data induk} titik pantau beserta nilai \textbf{terakhir}nya (satu baris per titik).
  \item \kode{sensor\_logs} menyimpan \textbf{seluruh riwayat} pengukuran (satu baris per pesan). Jika sebuah sensor dihapus, lognya ikut terhapus (\kode{ON DELETE CASCADE}).
  \item \kode{notifications} menyimpan peringatan. Jika sensor dihapus, notifikasi tetap ada tetapi \kode{sensor\_id} dikosongkan (\kode{ON DELETE SET NULL}).
\end{itemize}

\subsection{Kamus Data}

\begin{table}[H]
\centering
\caption{Kamus data tabel \kode{sensors}}
\label{tab:kd-sensors}
\begin{tabularx}{\linewidth}{@{}llL@{}}
\toprule
\textbf{Kolom} & \textbf{Tipe} & \textbf{Keterangan} \\
\midrule
id               & TEXT, PK         & Kode titik pantau, mis. \kode{PG-001}. \\
name             & TEXT             & Nama titik pantau. \\
region           & TEXT             & Wilayah / keterangan lokasi. \\
latitude         & DOUBLE           & Garis lintang. \\
longitude        & DOUBLE           & Garis bujur. \\
last\_moisture   & DOUBLE           & Kelembapan terakhir (\%), bawaan 65,0. \\
last\_ph         & DOUBLE           & pH terakhir, bawaan 4,2. \\
last\_temperature& DOUBLE           & Suhu terakhir (\textdegree C), bawaan 28,5. \\
status           & TEXT             & \kode{normal} / \kode{warning} / \kode{critical}. \\
updated\_at      & TIMESTAMPTZ      & Waktu pembaruan terakhir (UTC). \\
\bottomrule
\end{tabularx}
\end{table}

\begin{table}[H]
\centering
\caption{Kamus data tabel \kode{sensor\_logs}}
\label{tab:kd-logs}
\begin{tabularx}{\linewidth}{@{}llL@{}}
\toprule
\textbf{Kolom} & \textbf{Tipe} & \textbf{Keterangan} \\
\midrule
id          & BIGINT, PK   & Nomor urut otomatis. \\
created\_at & TIMESTAMPTZ  & Waktu data dicatat (UTC). Diberi indeks untuk pengurutan cepat. \\
sensor\_id  & TEXT, FK     & Mengacu ke \kode{sensors.id}. Diberi indeks. \\
sensor\_name& TEXT         & Nama titik saat data dicatat. \\
moisture    & DOUBLE       & Kelembapan tanah (\%). \\
ph          & DOUBLE       & pH tanah. \\
temperature & DOUBLE       & Suhu tanah (\textdegree C). \\
latitude, longitude & DOUBLE & Posisi saat pengukuran (boleh kosong). \\
status      & TEXT         & Wajib salah satu: \kode{normal}, \kode{warning}, \kode{critical}. \\
reason      & TEXT         & Kalimat penjelasan mengapa status tersebut diberikan. \\
\bottomrule
\end{tabularx}
\end{table}

\begin{table}[H]
\centering
\caption{Kamus data tabel \kode{notifications}}
\label{tab:kd-notif}
\begin{tabularx}{\linewidth}{@{}llL@{}}
\toprule
\textbf{Kolom} & \textbf{Tipe} & \textbf{Keterangan} \\
\midrule
id          & UUID, PK     & ID acak otomatis. \\
created\_at & TIMESTAMPTZ  & Waktu notifikasi dibuat. \\
sensor\_id  & TEXT, FK     & Titik pantau terkait (boleh kosong). \\
title       & TEXT         & Judul, mis. ``BAHAYA: Kondisi Kritis di Baleendah''. \\
message     & TEXT         & Isi pesan (alasan status). \\
severity    & TEXT         & \kode{info} / \kode{warning} / \kode{critical}. \\
is\_read    & BOOLEAN      & Sudah dibaca atau belum, bawaan \kode{false}. \\
\bottomrule
\end{tabularx}
\end{table}

\subsection{Keamanan Tingkat Baris (RLS) dan Realtime}
Ketiga tabel mengaktifkan \textit{Row Level Security} (RLS). Kebijakan yang dipasang saat ini:

\begin{table}[H]
\centering
\caption{Kebijakan akses (RLS) saat ini}
\label{tab:rls}
\begin{tabular}{@{}lcccc@{}}
\toprule
\textbf{Tabel} & \textbf{Baca} & \textbf{Tambah} & \textbf{Ubah} & \textbf{Hapus} \\
\midrule
sensors        & publik & publik & publik & publik \\
sensor\_logs   & publik & publik & --     & --     \\
notifications  & publik & publik & publik & --     \\
\bottomrule
\end{tabular}
\end{table}

Ketiga tabel juga didaftarkan ke publikasi \kode{supabase\_realtime}, sehingga setiap
perubahan data dapat langsung ``didengar'' oleh \textit{dashboard}.

% =====================================================================
\section{Aturan Penentuan Status Risiko}
% =====================================================================

Inti ``kecerdasan'' PeatGuard adalah aturan sederhana berbasis rujukan ilmiah
(Huang \& Rein 2015; Usup dkk. 2004; Saharjo \& Nurhayati 2006; PP No.~57 Tahun 2016).
Aturan yang sama dipakai di \textit{bridge worker} dan di \textit{dashboard}.

\subsection{Tabel Ambang Batas}
\begin{table}[H]
\centering
\caption{Ambang batas parameter sensor}
\label{tab:ambang}
\begin{tabularx}{\linewidth}{@{}lccL@{}}
\toprule
\textbf{Parameter} & \textbf{Aman} & \textbf{Waspada} & \textbf{Kritis} \\
\midrule
Kelembapan tanah (\%) & $\geq 60$ & $40$ -- $59{,}9$ & $< 40$ \\
Suhu tanah (\textdegree C) & $< 32$ & $32$ -- $34{,}9$ & $\geq 35$ \\
pH tanah & $3{,}2$ -- $5{,}5$ & $< 3{,}2$ atau $> 5{,}5$ & (tidak dipakai untuk status kritis) \\
\bottomrule
\end{tabularx}
\end{table}

\textbf{Mengapa angka-angka ini?}
\begin{itemize}
  \item \textbf{Kelembapan $<$ 40\%}: gambut sangat kering, bara api mudah menyala dan menjalar
        (batas \textit{Critical Moisture of Extinction}).
  \item \textbf{Suhu $\geq$ 35\textdegree C} di kedalaman sensor: tanda kuat adanya bara panas di bawah permukaan.
  \item \textbf{pH $<$ 3,2}: tanah terlalu kering sehingga mineral pirit teroksidasi.
        \textbf{pH $>$ 5,5}: kemungkinan ada abu sisa pembakaran.
\end{itemize}

\subsection{Diagram Alir Keputusan}
Urutan pemeriksaan penting: kondisi \textbf{kritis} diperiksa lebih dulu, baru \textbf{waspada}.

\begin{figure}[H]
\centering
\begin{tikzpicture}[node distance=0.7cm]
  \node[startnode] (s) {};
  \node[action, below=of s] (in) {Ambil nilai terbaru:\\ K (kelembapan), S (suhu), pH};
  \node[decision, below=of in] (d1) {K $<$ 40\\atau S $\geq$ 35 ?};
  \node[decision, below=1.0cm of d1] (d2) {K $<$ 60 atau S $\geq$ 32\\atau pH $<$ 3,2 atau pH $>$ 5,5 ?};
  \node[action, below=1.0cm of d2, minimum width=3cm, text width=2.8cm] (ok) {\textbf{AMAN}\\ (normal)};
  \node[action, right=1.4cm of d1, minimum width=3cm, text width=2.8cm] (kr) {\textbf{KRITIS}\\ (critical)};
  \node[action, right=0.6cm of d2, minimum width=3cm, text width=2.8cm] (ws) {\textbf{WASPADA}\\ (warning)};
  \node[endnode, below=of ok] (e) {};

  \draw[flow] (s) -- (in);
  \draw[flow] (in) -- (d1);
  \draw[flow] (d1) -- node[lbl, above]{ya} (kr);
  \draw[flow] (d1) -- node[lbl, right]{tidak} (d2);
  \draw[flow] (d2) -- node[lbl, above]{ya} (ws);
  \draw[flow] (d2) -- node[lbl, right]{tidak} (ok);
  \draw[flow] (ok) -- (e);
  \draw[flow] (kr.east) -- ++(0.5,0) |- (e.east);
  \draw[flow] (ws.east) -- ++(0.5,0) |- (e.east);
\end{tikzpicture}
\caption{Diagram alir penentuan status risiko}
\label{fig:flow-status}
\end{figure}

\subsection{Contoh Perhitungan}
\begin{table}[H]
\centering
\caption{Contoh hasil evaluasi status}
\label{tab:contoh}
\begin{tabular}{@{}cccll@{}}
\toprule
\textbf{K (\%)} & \textbf{S (\textdegree C)} & \textbf{pH} & \textbf{Status} & \textbf{Alasan utama} \\
\midrule
68,4 & 28,5 & 4,2 & Aman    & Semua dalam batas normal \\
41,8 & 33,7 & 3,8 & Waspada & Kelembapan di zona defisit (40--60\%) \\
65,0 & 28,0 & 5,9 & Waspada & pH $>$ 5,5 (indikasi abu) \\
38,5 & 30,0 & 4,0 & Kritis  & Kelembapan di bawah 40\% \\
62,0 & 36,1 & 4,3 & Kritis  & Suhu $\geq$ 35\textdegree C (indikasi bara) \\
\bottomrule
\end{tabular}
\end{table}

Potongan kode aturan tersebut (dari \kode{scripts/mqtt\_supabase\_bridge.js}):
\begin{lstlisting}
function evaluateStatus(m, t, p) {
  if (m < 40.0 || t >= 35.0) return "critical";
  if (m < 60.0 || t >= 32.0 || p < 3.2 || p > 5.5) return "warning";
  return "normal";
}
\end{lstlisting}

% =====================================================================
\section{Komunikasi MQTT}
% =====================================================================

\subsection{Daftar Topik}
Pola topik: \kode{sensor/soil/<parameter><nomor\_titik>}. \textit{Bridge} dan \textit{dashboard}
berlangganan ke \kode{sensor/soil/\#} (tanda \kode{\#} artinya ``semua topik di bawahnya'').

\begin{table}[H]
\centering
\caption{Topik MQTT yang digunakan}
\label{tab:topik}
\begin{tabular}{@{}lllc@{}}
\toprule
\textbf{Topik} & \textbf{Parameter} & \textbf{Titik} & \textbf{Satuan} \\
\midrule
\kode{sensor/soil/moisture1} & Kelembapan & Parongpong & \% \\
\kode{sensor/soil/ph1}       & pH         & Parongpong & pH \\
\kode{sensor/soil/temp1}     & Suhu       & Parongpong & \textdegree C \\
\kode{sensor/soil/moisture2} & Kelembapan & Baleendah  & \% \\
\kode{sensor/soil/ph2}       & pH         & Baleendah  & pH \\
\kode{sensor/soil/temp2}     & Suhu       & Baleendah  & \textdegree C \\
\kode{sensor/soil/moisture3} & Kelembapan & Rancaupas  & \% \\
\kode{sensor/soil/ph3}       & pH         & Rancaupas  & pH \\
\kode{sensor/soil/temp3}     & Suhu       & Rancaupas  & \textdegree C \\
\bottomrule
\end{tabular}
\end{table}

\subsection{Format Pesan}
Sistem menerima dua bentuk pesan:

\textbf{1. Angka biasa} (paling sederhana):
\begin{lstlisting}
33.7
\end{lstlisting}

\textbf{2. JSON} (dapat membawa lokasi dan waktu):
\begin{lstlisting}
{
  "value": 33.7,
  "latitude": -6.8320,
  "longitude": 107.5910,
  "timestamp": "2026-11-23T08:15:00Z",
  "sensor": "temp2"
}
\end{lstlisting}

Jika nama topik tidak mengikuti pola, sistem mencoba membaca jenis parameter dan nomor titik dari
kolom \kode{sensor}. Pesan yang tidak berisi angka atau tidak dikenali nomor titiknya akan diabaikan.

\begin{catatan}[Catatan penting]
Satu pesan hanya membawa \textbf{satu} parameter. Karena itu, \textit{bridge} menyimpan nilai terakhir
setiap parameter (\textit{cache}) agar status bisa dihitung dari kombinasi kelembapan, suhu, dan pH terbaru.
\end{catatan}

% =====================================================================
\section{Spesifikasi API Backend}
% =====================================================================

Alamat dasar: \kode{http://127.0.0.1:8000}. Dokumentasi interaktif otomatis tersedia di
\kode{/docs} (Swagger UI) setelah server berjalan.

\begin{table}[H]
\centering
\caption{Daftar endpoint API}
\label{tab:api}
\begin{tabularx}{\linewidth}{@{}llL@{}}
\toprule
\textbf{Metode} & \textbf{Endpoint} & \textbf{Fungsi} \\
\midrule
GET  & \kode{/}                                   & Cek status layanan. \\
GET  & \kode{/api/search?q=...}                   & Mencari sensor dan halaman berdasarkan kata kunci. Jika \kode{q} kosong, mengembalikan semua sensor sebagai saran. \\
GET  & \kode{/api/notifications}                  & Mengambil semua notifikasi beserta jumlah yang belum dibaca. \\
POST & \kode{/api/notifications}                  & Menambah notifikasi baru. \\
POST & \kode{/api/notifications/\{id\}/read}      & Menandai satu notifikasi sebagai dibaca. \\
POST & \kode{/api/notifications/mark-all-read}    & Menandai semua notifikasi sebagai dibaca. \\
\bottomrule
\end{tabularx}
\end{table}

\subsection{Contoh: Pencarian}
Permintaan:
\begin{lstlisting}
GET /api/search?q=waspada
\end{lstlisting}
Jawaban:
\begin{lstlisting}
{
  "query": "waspada",
  "total_results": 2,
  "results": [
    { "id": "PG-002", "title": "Titik Pantau 2 - Sungai Tohor",
      "category": "sensor", "status": "warning", "url": "/riwayat",
      "subtitle": "PG-002 · Tebing Tinggi Timur, Riau · 41.8% kelembapan · 3.9 pH" },
    { "id": "page-peringatan", "title": "Peringatan Dini Lahan",
      "category": "halaman", "status": null, "url": "/peringatan",
      "subtitle": "Daftar titik pantau dengan kondisi kritis atau waspada kebakaran" }
  ]
}
\end{lstlisting}

Kata kunci khusus yang dipahami pencarian: \kode{aman} (status normal), \kode{waspada} (warning),
\kode{bahaya} (critical), \kode{kering} (semua yang tidak normal).

\subsection{Contoh: Menambah Notifikasi}
Permintaan:
\begin{lstlisting}
POST /api/notifications
Content-Type: application/json

{
  "title": "Kondisi Gambut Kering (Waspada)",
  "message": "Titik Pantau 2 mencatat kelembapan 41.8%.",
  "type": "warning",
  "sensor_id": "PG-002",
  "url": "/peringatan"
}
\end{lstlisting}

Jawaban berupa objek notifikasi lengkap dengan \kode{id}, \kode{created\_at}, dan \kode{is\_read = false}.
Nilai \kode{type} yang dikenal: \kode{critical}, \kode{warning}, \kode{info}, \kode{success}.

% =====================================================================
\section{Halaman Dashboard}
% =====================================================================

\begin{table}[H]
\centering
\caption{Daftar halaman dashboard}
\label{tab:halaman}
\begin{tabularx}{\linewidth}{@{}llL@{}}
\toprule
\textbf{Menu} & \textbf{Alamat} & \textbf{Isi utama} \\
\midrule
Dashboard      & \kode{/}               & Ringkasan angka tiap titik, status koneksi MQTT dan Supabase, 9 topik telemetri, peta mini. \\
Riwayat Data   & \kode{/riwayat}        & Grafik tren (kelembapan/pH/suhu) dan tabel log yang dapat dicari, diurutkan, dan dipaginasi. \\
Peta Sensor    & \kode{/peta-sensor}    & Peta Leaflet dengan penanda titik pantau berdasarkan status. \\
Kondisi Lahan  & \kode{/kondisi-lahan}  & Rata-rata kondisi fisik \& kimia, tabel rincian per titik, kartu rujukan ambang batas. \\
Peringatan     & \kode{/peringatan}     & Daftar titik berstatus Waspada/Kritis lengkap dengan alasan. \\
Laporan        & \kode{/laporan}        & Kesimpulan lapangan, daftar titik, dan tombol ekspor CSV. \\
\bottomrule
\end{tabularx}
\end{table}

\subsection{Sumber Data Riwayat (Urutan Prioritas)}
Saat \textit{dashboard} dibuka, data riwayat dimuat dengan urutan berikut:
\begin{enumerate}
  \item \textbf{Supabase} -- 150 log terbaru. Jika berhasil, status Supabase menjadi ``terhubung''.
  \item \textbf{Penyimpanan lokal browser} (\kode{localStorage}) -- jika Supabase gagal atau kosong.
  \item \textbf{Data contoh} -- dibuat otomatis (37 titik waktu tiap 40 menit untuk 3 sensor) jika keduanya kosong.
\end{enumerate}
Di \textit{browser}, maksimal 200 log terbaru dipertahankan.

% =====================================================================
\section{Cara Menjalankan Sistem}
% =====================================================================

\subsection{Backend (be-peat)}
\begin{lstlisting}
cd be-peat
pip install -r requirements.txt
python run.py
# Server berjalan di http://127.0.0.1:8000  (dokumentasi: /docs)
\end{lstlisting}

\subsection{Basis Data (Supabase)}
Buka \textit{SQL Editor} di proyek Supabase, lalu jalankan isi berkas
\kode{dashboard-peat/supabase/schema.sql}. Skrip ini aman dijalankan berulang kali (idempoten):
membuat tabel, indeks, kebijakan RLS, data awal 3 titik, dan mengaktifkan \textit{Realtime}.

\subsection{Dashboard dan Bridge (dashboard-peat)}
\begin{lstlisting}
cd dashboard-peat
npm install
npm run dev       # mode pengembangan -> http://localhost:5173
npm run bridge    # jalankan bridge MQTT -> Supabase (terminal terpisah)
\end{lstlisting}

\subsection{Produksi dengan Docker}
\begin{lstlisting}
docker build -t peatguard-dashboard .
docker run -p 3000:3000 peatguard-dashboard
\end{lstlisting}
\textit{Bridge worker} tidak termasuk di dalam kontainer tersebut, sehingga perlu dijalankan sebagai
proses/layanan terpisah (misalnya dengan \kode{pm2}, \kode{systemd}, atau kontainer kedua).

% =====================================================================
\section{Temuan Teknis dan Rekomendasi}
\label{sec:temuan}
% =====================================================================

Bagian ini merangkum hal-hal yang perlu diperbaiki sebelum sistem dipakai secara luas.
Tingkat prioritas: \textbf{T} = tinggi, \textbf{S} = sedang, \textbf{R} = rendah.

\begin{longtable}{@{}p{0.6cm}p{4.6cm}p{8.6cm}@{}}
\caption{Temuan dan rekomendasi}\label{tab:temuan}\\
\toprule
\textbf{Prio.} & \textbf{Temuan} & \textbf{Rekomendasi} \\
\midrule
\endfirsthead
\toprule
\textbf{Prio.} & \textbf{Temuan} & \textbf{Rekomendasi} \\
\midrule
\endhead
\bottomrule
\endfoot
T & Kata sandi MQTT dan kunci Supabase ditulis langsung di kode (\kode{SensorContext.jsx}, \kode{mqtt\_supabase\_bridge.js}, \kode{lib/supabase.js}) dan ikut terkirim ke \textit{browser}. &
Pindahkan ke variabel lingkungan (\kode{.env}); ganti kata sandi MQTT yang sudah terlanjur tersimpan di riwayat Git; buat akun MQTT khusus \textit{dashboard} yang \textbf{hanya boleh membaca} (\textit{subscribe}). \\
T & Kebijakan RLS mengizinkan siapa pun menambah, mengubah, bahkan menghapus data \kode{sensors} serta menambah log dan notifikasi. &
Batasi penulisan hanya untuk \textit{bridge} memakai kunci \textit{service role} di sisi server; publik cukup hak baca. \\
T & Data log tercatat ganda: \textit{bridge} \textbf{dan} setiap \textit{browser} yang terbuka sama-sama menyimpan ke \kode{sensor\_logs}. &
Jadikan \textit{bridge} satu-satunya penulis. \textit{Dashboard} cukup membaca dan mendengarkan \textit{Realtime}. \\
S & \textit{Dashboard} memanggil API di \kode{http://127.0.0.1:8000}, dan \textit{backend} hanya mendengarkan di \kode{127.0.0.1}. Ini hanya berfungsi bila \textit{browser} dan \textit{backend} berada di komputer yang sama. &
Gunakan alamat API yang dapat dikonfigurasi (mis. \kode{VITE\_API\_URL}) dan jalankan Uvicorn di \kode{0.0.0.0} di balik \textit{reverse proxy} HTTPS. \\
S & Notifikasi di \textit{backend} disimpan di memori: hilang saat server dimatikan dan tidak tersambung dengan tabel \kode{notifications} di Supabase. &
Hubungkan \kode{NotificationController} ke tabel \kode{notifications} Supabase agar sumber notifikasi hanya satu. \\
S & Data sensor di \textit{backend} (nama dan wilayah di Riau) berbeda dengan \textit{dashboard} dan basis data (Parongpong, Baleendah, Rancaupas -- Bandung). Data di \textit{backend} juga statis, bukan nilai \textit{realtime}. &
Ambil data sensor untuk pencarian langsung dari tabel \kode{sensors}. \\
S & CORS \textit{backend} mengizinkan semua asal (\kode{*}) sekaligus \kode{allow\_credentials}. &
Batasi ke domain \textit{dashboard} yang resmi. \\
R & Nilai \kode{type}/\kode{severity} notifikasi tidak seragam: \textit{backend} memakai \kode{success}, sedangkan basis data hanya menerima \kode{info}, \kode{warning}, \kode{critical}. &
Samakan daftar nilai di semua bagian. \\
R & Aturan status ditulis ulang di dua tempat (\kode{data/sensors.js} dan \textit{bridge}); dokumen rujukan menyebut ambang pH kritis ($<2{,}8$ / $>6{,}5$) yang belum diterapkan di kode. &
Satukan aturan dalam satu modul bersama dan putuskan apakah ambang pH kritis akan dipakai. \\
R & Notifikasi di Supabase dibuat \textbf{setiap} pesan berstatus kritis, sehingga bisa membanjiri tabel. &
Buat notifikasi hanya saat status \textbf{berubah} menjadi kritis, atau beri jeda waktu minimum. \\
R & Belum ada pengujian otomatis (\textit{unit test}) maupun autentikasi pengguna \textit{dashboard}. &
Tambahkan \textit{test} untuk fungsi evaluasi status dan API; tambahkan login (mis. Supabase Auth) bila data bersifat internal. \\
\end{longtable}

% =====================================================================
\section{Penutup}
% =====================================================================

PeatGuard menggabungkan sensor IoT, protokol MQTT, basis data \textit{cloud} dengan fitur \textit{realtime},
dan \textit{dashboard} web untuk memberikan \textbf{peringatan dini kebakaran lahan gambut}.
Logika utamanya sederhana dan transparan: tiga parameter tanah dibandingkan dengan ambang batas ilmiah
untuk menghasilkan status Aman, Waspada, atau Kritis.

Secara fungsional sistem sudah berjalan dari ujung ke ujung. Langkah berikutnya yang disarankan adalah
memperkuat sisi keamanan (kredensial dan hak akses basis data), menghilangkan pencatatan ganda,
serta menyatukan sumber data notifikasi dan sensor agar sistem siap digunakan di lapangan secara luas.

\subsection*{Rujukan}
\begin{enumerate}[label={[\arabic*]}]
  \item Huang, X., \& Rein, G. (2015). Smouldering combustion of peat in wildfires. \textit{International Journal of Wildland Fire}.
  \item Usup, A., Hashimoto, Y., Takahashi, H., \& Hayasaka, H. (2004). Combustion characteristics of peat in Central Kalimantan, Indonesia. \textit{Tropics}, 14(1), 1--19.
  \item Saharjo, B. H., \& Nurhayati, A. D. (2006). Karakteristik Pembakaran dan Suhu Bawah Permukaan Lahan Gambut Terdegradasi. \textit{Jurnal Silvikultur Tropika}, IPB.
  \item Page, S. E., Rieley, J. O., \& Banks, C. J. (2011). Global and regional importance of the tropical peatland carbon pool. \textit{Global Change Biology}, 17(2), 798--818.
  \item Peraturan Pemerintah RI No.~57 Tahun 2016 tentang Perlindungan dan Pengelolaan Ekosistem Gambut.
  \item Dokumentasi resmi: FastAPI, React Router, Supabase, MQTT.js, Leaflet.
\end{enumerate}

\end{document}
